\documentclass{beamer}





\usepackage[utf8]{inputenc}
\usepackage[german]{babel}
\usepackage{hyperref}
\usepackage{chronosys}
\usepackage[utf8]{inputenc}

\usetheme{Montpellier}
\usecolortheme{default}
\setbeamercolor{background canvas}{bg=white}
\setbeamercolor{normal text}{fg=black}

\title{Geschichte der Fotografie}
\author{Jurek Aita} % hier noch eure name eintragen
\date{\today}

\begin{document}

\frame{\titlepage}

\begin{frame}
\frametitle{Inhaltsverzeichnis}
\tableofcontents
\end{frame}

\section{Die Camera Obscura}
\begin{frame}
    \begin{itemize}
        \item Erster bekannte Vorläufer der Camera
        \item 
    \end{itemize}
\end{frame}
\section{Camera Obscura Funktionsweiße}
\begin{frame}
    \href{https://bridge.klett.de/EWA-RI9ITG79N5/media/interactive/3423291/index.html}{Hier zur Simulation}
\end{frame}

\section{Geschichte Camera Obscura}
\begin{frame}
    \frametitle{Projekt-Timeline}
    \centering
    % Der Zeitstrahl wird automatisch auf die Textbreite skaliert
    \only<1>{
        \startchronology[startyear=-300, stopyear=1700, height=2ex]
            \chronoevent{-300}{}
        \stopchronology
    }
    \only<2>{
        \startchronology[startyear=-300, stopyear=1700, height=2ex]
            \chronoevent{980}{}
        \stopchronology
    }
    \only<3>{
        \startchronology[startyear=-300, stopyear=1700, height=2ex]
            \chronoevent{1300}{}
        \stopchronology
    }
    \only<4>{
        \startchronology[startyear=-300, stopyear=1700, height=2ex]
            \chronoevent{1400}{}
        \stopchronology
    }
    \only<5>{
        \startchronology[startyear=-300, stopyear=1700, height=2ex]
            \chronoevent{1600}{}
        \stopchronology
    }
    
    \vspace{1.5cm}
    
    \only<1>{\textbf{\Large Beschreibung der physikalischen Eigenschaften}}
    \only<2>{\textbf{\Large Erster bekannter Versuch}}
    \only<3>{\textbf{\Large Praktische Anwendung}}
    \only<4>{\textbf{\Large Leonardo da Vinci}}
    \only<5>{\textbf{\Large Mobile Version}}
\end{frame}

\section{Camera Obscura in der Kunst}

\section{Massenmedium Fotografie}
\section{Portraitfotografie und Inzinierung des 19. Jahrhunderts}
\section{Selbstbild vs. Fremdbild}

\end{document}